%\begin{center}
%\large \bf \runtitulo
%\end{center}
%\vspace{1cm}
\chapter*{Proving the execution of a high level language with Plonky2}

This work studies the possibility of proving the execution of high-level programs using Plonky2, a ZK-SNARK cryptographic proving system that does not require a \textit{trusted setup}. The work is situated within the field of Zero-Knowledge Proofs, presenting an introduction to the mathematical and cryptographic primitives necessary to understand these protocols, as well as a high-level analysis of relevant systems such as PLONK.


A tool is proposed and developed that translates ACIR programs (the intermediate representation of Noir) into Plonky2. In this way, the construction of verifiable execution proofs in environments with limited storage resources is achieved, overcoming one of the main limitations of Barretenberg.


The document includes a formal description of the translation of primitives and operations, an experimental evaluation of the achieved performance, and a comparative analysis in terms of execution times and the sizes of the generated artifacts. Finally, the results obtained are discussed and directions for future work are outlined.


\bigskip

\noindent\textbf{Keywords:} Zero Knowledge Proofs, ZK, SNARK, PLONK, Plonky2, Noir, ACIR, Barretenberg, Criptografía.